% capitulos/cap5/tema.tex -- §5.9: tematização dinâmica (paletas, neutros por modo, contraste, leitura da configuração).
% Parte do Capítulo 5 (capitulos/cap5_arquitetura.tex); os rótulos \label e os números de seção não mudaram.

\section{Tematização Dinâmica}\label{sec:arq-tema}

A personalização visual (RF04, RF18, RF19, RF31, RF32) baseia-se em variáveis CSS alimentadas pelo banco. Esta seção descreve primeiro o mecanismo \emph{atual} (administrador único, tema único e claro), que está integrado, e depois, nas subseções seguintes, o \emph{alvo} da \adr{ADR-025}, da \adr{ADR-030} e da \adr{ADR-031} (aparência por proprietário, duas paletas e modo noturno; módulos M15 e M18). O mecanismo atual, ilustrado na Figura~\ref{fig:seq-personalizacao-tema}, tem três partes:

\begin{enumerate}
  \item \textbf{Tokens.} Em \arq{globals.css}, o bloco \texttt{@theme inline} do Tailwind mapeia os tokens de cor (por exemplo, \arq{--color-primary}) para variáveis (por exemplo, \arq{var(--vitrine-primary)}); os utilitários compilados referenciam a variável, não o valor;
  \item \textbf{Injeção.} O \en{layout} raiz é estático e não lê o banco (merge M-021). O portal, o \texttt{/login} e o painel, Server Components, leem a configuração do site e a definem no atributo \texttt{style} do envoltório \texttt{.vitrine-tema} (\arq{src/app/_lib/site-theme.js}) --- as duas paletas (\arq{--vitrine-light-*} e \arq{--vitrine-dark-*}) e o texto sobre a primária e a secundária de cada modo, calculado por \texttt{textColorsOf} (Seção~\ref{sec:arq-contraste}) ---, além do nome e da descrição do site nos metadados do próprio layout (o logotipo e o \en{favicon} entram quando houver componente que os exiba); os valores são validados como \texttt{\#RRGGBB} antes de persistir e, por desconfiança do dado lido, de novo na leitura, o que impede a injeção de CSS (Seção~\ref{sec:arq-tema-estatico});
  \item \textbf{Invalidação.} Ao salvar a configuração, o Route Handler invalida a página do portal (\texttt{revalidatePath("/")}, literal; nunca o tipo \texttt{layout}); na próxima requisição, ela é regenerada com as novas variáveis, sem \en{build}, e o login e o painel, dinâmicos, já nascem com a configuração (Seção~\ref{sec:arq-tema-estatico}).
\end{enumerate}

\begin{lstlisting}[language=VCSS, caption={Tokens de tema em globals.css (resumo; merges M-005 e M-009)}, label={lst:tema-css}]
@import "tailwindcss";

:root {
  /* Paleta configurável (contrato do M11): <html style> sobrescreve o padrão */
  --vitrine-primary: #3b82f6;
  --vitrine-secondary: #1e293b;
  --vitrine-accent: #f59e0b;
  /* Texto sobre a primária = readableTextOn(primária): branco ou preto; definida no layout raiz */
  --vitrine-on-primary: #000000;
  --vitrine-on-secondary: #ffffff;

  /* Neutros fixos (modo claro) e foco independente */
  --vitrine-background: #ffffff;  --vitrine-foreground: #171717;
  --vitrine-surface: #f5f5f5;     --vitrine-muted: #525252;
  --vitrine-border: #d4d4d4;      --vitrine-border-strong: #737373;
  --vitrine-danger: #b91c1c;
  --vitrine-focus: #171717;       --vitrine-focus-halo: #ffffff;

  /* Links do conteúdo (M08): a cor do texto, distinguidos pelo sublinhado */
  --vitrine-link: var(--vitrine-foreground);
  /* Realce de código (M08), tema claro, com 4,5:1 sobre o fundo do bloco: além de
     background e foreground, comment, keyword, string, number, function, type, variable e deletion */
  --vitrine-code-background: #f5f5f5;  --vitrine-code-foreground: #171717;
  --vitrine-code-keyword: #be123c;     --vitrine-code-string: #047857;
}

/* ADR-025: o modo noturno vale sob a classe .dark; por ora, só o realce de código o define */
.dark {
  --vitrine-code-background: #171717;  --vitrine-code-foreground: #f5f5f5;
  --vitrine-code-keyword: #fda4af;     --vitrine-code-string: #6ee7b7;
}

@theme inline {
  --color-primary: var(--vitrine-primary);
  --color-secondary: var(--vitrine-secondary);
  --color-accent: var(--vitrine-accent);
  --color-on-primary: var(--vitrine-on-primary);
  --color-on-secondary: var(--vitrine-on-secondary);
  /* neutros no mesmo padrão: --color-background, -foreground, -surface... */
}

/* a variante dark: só vale sob a classe .dark, nunca pela preferência do sistema */
@custom-variant dark (&:where(.dark, .dark *));
\end{lstlisting}

O mapeamento semântico é: \textbf{primária} para ações e ênfase (botão principal e item ativo da navegação, com o fundo \texttt{bg-primary} e o texto \texttt{text-on-primary}); \textbf{secundária} para a estrutura (fundo da barra superior e cabeçalhos, com texto branco, \texttt{text-on-secondary}) e para o \en{hover} dos botões; e \textbf{destaque} para marcadores e realce não textual (marcador do item ativo e \en{hover} da navegação). O \textbf{indicador de foco} não usa a cor de destaque: é neutro e global --- contorno de 2~px em \texttt{\#171717} com halo branco (\arq{--vitrine-focus} e \arq{--vitrine-focus-halo}) ---, de modo que fica visível sobre qualquer fundo, inclusive sobre uma paleta que o administrador troque (RNF05; \adr{ADR-020}). Os demais tons da interface --- fundo, texto, superfícies, bordas e mensagens de erro --- são neutros fixos, não configuráveis. A \adr{ADR-025} substitui o \decisao{D-17} (sem modo escuro): o modo noturno é especificado na Seção~\ref{sec:arq-tema-padrao}. Até o M18 estar integrado, vale somente o modo claro, com a paleta escolhida pelo administrador, e a variante \texttt{dark:} do Tailwind só vale sob a classe \texttt{.dark}, nunca pela preferência do sistema; o merge M-009 já define, sob \texttt{.dark}, as cores noturnas do realce de código. Por isso, hoje, o fundo da página é fixo e branco, o texto sobre a cor secundária é branco, o texto sobre a cor primária é escolhido por \texttt{readableTextOn} (Seção~\ref{sec:arq-contraste}) e só as três cores do tema são verificadas. Os componentes usam apenas os utilitários gerados pelos tokens (por exemplo, \texttt{bg-primary}, \texttt{text-on-primary}, \texttt{bg-surface} e \texttt{text-muted}), nunca as cores fixas do Tailwind (RF19).

\subsection{Paletas Padrão, Neutros por Modo e Alternância}\label{sec:arq-tema-padrao}

Esta seção é a especificação do módulo M18 (\adr{ADR-025}, \adr{ADR-031}); a etapa 1 da UI foi integrada no merge M-017 (05/10/2026): os neutros noturnos, as paletas padrão da \adr{ADR-031}, o envoltório \texttt{.vitrine-tema}, o contraste por modo (\arq{theme.js}), o \en{script} do modo e o seletor já existem como componentes e módulos isolados; o formulário de aparência e a ligação nas páginas (layout raiz e layout do espaço) são das etapas seguintes e da Parte 3 da A-014, e, até o Core mover a paleta para o \texttt{.vitrine-tema}, o \texttt{style} do \texttt{<html>} do layout atual ainda fixa a paleta clara no modo noturno. Cada proprietário tem duas paletas de três cores (primária, secundária e destaque) sobre \emph{neutros fixos} de cada modo; só as paletas são configuráveis, o que mantém o contraste verificável com poucos pares por modo (RN27).

\textbf{Paletas padrão.} Os valores da Tabela~\ref{tab:padroes-cor} são os padrões do esquema, da leitura de uma aparência ainda não gravada e do \en{seed}; nenhum gera aviso de contraste. Em uma instalação existente, a migração do M15 trocou, na linha do site, só as cores que ainda valem o padrão antigo (\adr{ADR-031}).

\begin{table}[htb]
\caption{Paletas padrão e as razões de contraste com o fundo de cada modo}\label{tab:padroes-cor}
\small
\begin{tabularx}{\linewidth}{@{}L{2.2cm}L{2.0cm}L{2.0cm}L{2.0cm}Y@{}}
\toprule
\textbf{Cor} & \textbf{Modo claro} & \textbf{Razão com \texttt{\#FFFFFF}} & \textbf{Modo noturno} & \textbf{Razão com \texttt{\#0A0A0A}}\\
\midrule
Primária & \texttt{\#2563EB} & 5,17:1 (mín. 4,5) & \texttt{\#60A5FA} & 7,79:1 (mín. 4,5)\\
Secundária & \texttt{\#1E293B} & 14,63:1 (mín. 4,5) & \texttt{\#CBD5E1} & 13,33:1 (mín. 4,5)\\
Destaque & \texttt{\#B45309} & 5,02:1 (mín. 3) & \texttt{\#FBBF24} & 11,86:1 (mín. 3)\\
Texto sobre a primária & \texttt{\#FFFFFF} & 5,17:1 & \texttt{\#000000} & 8,26:1\\
Texto sobre a secundária & \texttt{\#FFFFFF} & 14,63:1 & \texttt{\#000000} & 14,14:1\\
\bottomrule
\end{tabularx}
\fonte{Elaborado pelos autores; razões calculadas com \texttt{contrastRatio} de \arq{src/lib/theme.js} e arredondadas (o aviso do sistema trunca, Seção~\ref{sec:arq-contraste}). O texto sobre cada cor é o resultado de \texttt{readableTextOn} e acompanha a cor, quando ela é personalizada.}
\end{table}

\textbf{Neutros por modo.} São definidos em \arq{globals.css} e nunca variam por proprietário (Tabela~\ref{tab:neutros}). O modo noturno vale sob a classe \texttt{.dark} em \texttt{<html>}, nunca pela preferência do sistema diretamente, para que a escolha do leitor prevaleça.

\begin{table}[htb]
\caption{Neutros fixos dos dois modos (razões com o fundo do próprio modo)}\label{tab:neutros}
\small
\begin{tabularx}{\linewidth}{@{}L{3.1cm}L{2.0cm}L{2.0cm}Y@{}}
\toprule
\textbf{Token} & \textbf{Claro} & \textbf{Noturno} & \textbf{Razões (claro; noturno) e uso}\\
\midrule
\arq{--vitrine-background} & \texttt{\#FFFFFF} & \texttt{\#0A0A0A} & Fundo da página\\
\arq{--vitrine-foreground} & \texttt{\#171717} & \texttt{\#EDEDED} & 17,92:1; 16,91:1. Texto principal\\
\arq{--vitrine-surface} & \texttt{\#F5F5F5} & \texttt{\#171717} & Cartões e fundo de código; o texto principal tem 16,44:1 e 15,31:1 sobre ele\\
\arq{--vitrine-muted} & \texttt{\#525252} & \texttt{\#A3A3A3} & 7,81:1; 7,84:1. Texto secundário (7,16:1 e 7,10:1 sobre a superfície)\\
\arq{--vitrine-border} & \texttt{\#D4D4D4} & \texttt{\#404040} & Divisores decorativos (sem exigência de contraste)\\
\arq{--vitrine-border-strong} & \texttt{\#737373} & \texttt{\#737373} & 4,74:1; 4,17:1. Borda de campos e controles (mínimo 3:1; 4,34:1 e 3,78:1 sobre a superfície)\\
\arq{--vitrine-danger} & \texttt{\#B91C1C} & \texttt{\#F87171} & 6,47:1; 7,15:1. Mensagens de erro (5,93:1 e 6,48:1 sobre a superfície)\\
Foco & \texttt{\#171717} & \texttt{\#FAFAFA} & Contorno de 2~px com halo do fundo (\arq{--vitrine-focus-halo}); 17,92:1 e 18,96:1\\
\bottomrule
\end{tabularx}
\fonte{Elaborado pelos autores; razões truncadas em duas casas, calculadas com \texttt{contrastRatio}.}
\end{table}

\textbf{Realce de código e KaTeX.} Os tokens \arq{--vitrine-code-*} e \arq{--vitrine-link} já têm variante noturna (merge M-009, com pior caso de 7,11:1). O KaTeX herda a cor do texto (\texttt{currentColor}) e dispensa variante própria; o M18 confirma isso nos dois modos.

\textbf{Aplicação das paletas.} O layout raiz é único e \emph{estático} (\adr{ADR-030}): não lê o banco e não define cores. O proprietário aparece em um elemento envoltório com a classe \texttt{vitrine-tema}, que o layout de \texttt{/\{handle\}} (aparência do usuário) e os do portal, do login e do painel (aparência do site) preenchem, no atributo \texttt{style}, com as variáveis \emph{das duas paletas}: \arq{--vitrine-light-primary}, \arq{--vitrine-light-secondary}, \arq{--vitrine-light-accent}, \arq{--vitrine-light-on-primary}, \arq{--vitrine-light-on-secondary} e as cinco equivalentes \arq{--vitrine-dark-*}. As duas regras do Código~\ref{lst:tema-modos} escolhem o conjunto do modo. Elas ficam no \emph{mesmo elemento} do atributo \texttt{style} de propósito: um \texttt{var()} declarado em \texttt{:root} seria resolvido ali, onde as variáveis do proprietário não existem. Os valores continuam validados como \texttt{\#RRGGBB} antes de persistir e na leitura (Seção~\ref{sec:arq-tema-estatico}), e os textos sobre as cores (\arq{on-*}) são calculados \emph{por modo}, com \texttt{readableTextOn}, porque o texto sobre a secundária deixa de ser sempre branco: no modo noturno com a paleta padrão é preto. A pré-visualização do formulário de aparência (M-023) usa duas classes auxiliares do \arq{globals.css}, \texttt{.vitrine-modo-claro} e \texttt{.vitrine-modo-escuro}, que forçam um modo seja qual for o da página, e as funções \texttt{DEFAULT\_PALETTES} e \texttt{themeStyleOf} de \arq{theme.js}.

\begin{lstlisting}[language=VCSS, caption={Escolha do conjunto de cores pelo modo (esboço normativo; vale também para as demais cores)}, label={lst:tema-modos}]
.vitrine-tema {
  --vitrine-primary: var(--vitrine-light-primary, #2563eb);
  --vitrine-on-primary: var(--vitrine-light-on-primary, #ffffff);
  /* secondary, accent e on-secondary, do mesmo modo */
}
.dark .vitrine-tema {
  --vitrine-primary: var(--vitrine-dark-primary, #60a5fa);
  --vitrine-on-primary: var(--vitrine-dark-on-primary, #000000);
}
\end{lstlisting}

\textbf{Alternância e primeira pintura.} O leitor escolhe Claro, Escuro ou Automático (RF32); a escolha fica em \texttt{localStorage}, na chave \texttt{vitrine:modo} (\texttt{light}, \texttt{dark} ou \texttt{auto}, que também é o valor de quem nunca escolheu). Um \en{script} inline no \texttt{<head>} do layout raiz (Código~\ref{lst:modo-cor}), de conteúdo \emph{fixo} (nunca montado com dado do banco ou do usuário), aplica a classe \texttt{.dark} antes da primeira pintura; a página não depende de \en{cookie} lido no servidor e permanece estática e com ISR (CT47). O \texttt{<html>} leva \texttt{suppressHydrationWarning}, porque a classe muda antes da hidratação. O \en{script} é admitido pela CSP (\texttt{script-src 'unsafe-inline'}, \adr{ADR-022}); o M18 o verifica sob \texttt{npm run build}, com o console aberto. O seletor de modo é um Componente Cliente que grava a chave e alterna a classe, e acompanha a mudança da preferência do sistema enquanto o modo é Automático. O CSS declara \texttt{color-scheme: light} em \texttt{:root} e \texttt{color-scheme: dark} em \texttt{.dark}, para que barras de rolagem e controles nativos acompanhem o modo.

\begin{lstlisting}[language=VJS, caption={Script do modo de cor no <head> do layout raiz (esboço normativo)}, label={lst:modo-cor}]
(function () {
  try {
    var m = localStorage.getItem("vitrine:modo");
    var escuro = m === "dark" ||
      (m !== "light" && matchMedia("(prefers-color-scheme: dark)").matches);
    if (escuro) document.documentElement.classList.add("dark");
  } catch (e) { /* armazenamento bloqueado: segue o modo claro */ }
})();
\end{lstlisting}

\textbf{Alvo (épico v1; \adr{ADR-036}, RF32, RN32): o modo por \en{cookie}, sem ``Automático''.} Quando o M22 for integrado, o parágrafo e o código acima são substituídos por isto: o leitor alterna entre \emph{dois} estados, Claro e Escuro; a escolha fica no \en{cookie} \texttt{vitrine-modo} (\texttt{light} ou \texttt{dark}; \texttt{Path=/}, um ano, \texttt{SameSite=Lax}, sem \texttt{HttpOnly}, \texttt{Secure} sob HTTPS), que o \emph{navegador} escreve (o seletor, ao clicar) e que só o \en{script} inline do \texttt{<head>} lê (Código~\ref{lst:modo-cor-v1}). \textbf{O servidor nunca lê esse \en{cookie}}: ler um \en{cookie} no layout o tornaria dinâmico e desligaria a ISR e o LCP; as páginas continuam estáticas (CT47). Sem o \en{cookie} (primeira visita), o modo é o \emph{claro}, mesmo que o sistema do leitor prefira o escuro; o \texttt{localStorage} \texttt{vitrine:modo} e o estado \texttt{auto} deixam de existir (não há migração: o valor antigo é ignorado). O \en{cookie} é funcional e não carrega dado pessoal (Seção~\ref{sec:arq-upload}). O mesmo \texttt{ColorModeToggle} fica nas páginas públicas e no cabeçalho do painel, e usa um \emph{par de ícones} em SVG \en{inline} do projeto, escolhido pelo administrador para o site inteiro entre quatro (\texttt{sun\_moon}, o padrão, \texttt{sun\_moon\_outline}, \texttt{lightbulb} e \texttt{contrast\_half}; coluna \texttt{modeIconPair} de \texttt{configurations}), recebido por \emph{prop}; o botão tem nome acessível e estado (``Mudar para o modo escuro'' ou ``claro'') e passa na varredura de acessibilidade.

\begin{lstlisting}[language=VJS, caption={Script do modo de cor no <head> com o cookie (alvo, esboço normativo)}, label={lst:modo-cor-v1}]
(function () {
  try {
    var m = document.cookie.match(/(?:^|; )vitrine-modo=(light|dark)/);
    if (m && m[1] === "dark") document.documentElement.classList.add("dark");
  } catch (e) { /* cookie bloqueado: segue o modo claro */ }
})();
\end{lstlisting}

\textbf{Cabeçalho do espaço e botão de modo (M-045; P-024).} A barra superior do espaço (\texttt{Navbar}) é \emph{um só cabeçalho}: o bloco do nome, da ``Administração'' e do botão de modo e o bloco das coleções têm o mesmo fundo (a cor secundária do modo), sem linha entre eles e sem a faixa de destaque que havia embaixo. Quem separa o cabeçalho do conteúdo é uma \textbf{sombra} (classe \texttt{.vitrine-cabecalho}, em \arq{globals.css}): suave no modo claro (a do \texttt{shadow-sm}) e mais forte no noturno (\texttt{0 2px 12px} preto a 40\%); como a pré-visualização da aparência força o modo pela classe do próprio envoltório, a regra vale também nos modos forçados da amostra. A sombra é preto translúcido, e não uma cor de tema, e é só decorativa: o contraste do texto continua medido sobre a cor secundária (Tabela~\ref{tab:pares-contraste}). O cabeçalho não é fixo (\emph{sticky}); isso é \emph{backlog} da v1.1. O botão do modo de cor traz o \emph{próprio} fundo neutro (\texttt{bg-background}, texto e borda neutros, 40$\times$40~px), de modo que o contraste dele não depende da paleta da barra: o fundo não aparece em volta do botão. Na pré-visualização de \arq{/admin/aparencia}, o botão aparece só como imagem (\texttt{ColorModeToggleStatic}, Seção~\ref{sec:arq-rotas}), com o par de ícones do site.

\textbf{Barra de coleções (M-046; P-026).} A faixa de coleções da \texttt{Navbar} segue a estética de uma barra de abas do projeto de referência do Engenheiro Chefe: os itens ficam encostados, com cantos retos e sem espaço em volta, até a base do cabeçalho. Em repouso o texto é \texttt{on-secondary} a 70\%; sob \emph{hover} ou foco e no item da coleção corrente o item se enche com a \textbf{cor primária do dono} (\texttt{bg-primary text-on-primary}). Só o item corrente tem, além disso, \texttt{aria-current}, o negrito e o sublinhado grosso. \textbf{Limite conhecido (decisão do Engenheiro Chefe, 11/10/2026):} o enchimento contra a barra tem contraste baixo com as paletas padrão (cerca de 2,8:1 no claro e 1,7:1 no noturno), e, se a primária do dono for igual à secundária, o enchimento desaparece; o item corrente continua identificável pelo negrito, pelo sublinhado e por \texttt{aria-current}, mas o \emph{hover} fica sem pista visível nesse caso. O aviso de contraste da aparência valida o texto sobre a primária, e não a primária contra a secundária (\emph{backlog} da v1.1). As setas que rolam a faixa são pílulas translúcidas derivadas do \texttt{on-secondary} (borda a 30\%, fundo a 10\%; pressionada, a primária), cujo ícone herda o contraste já validado do par \texttt{on-secondary} sobre \texttt{secondary}; a barra de rolagem fica oculta (\texttt{.vitrine-sem-barra}).

\subsection{Avisos de Contraste}\label{sec:arq-contraste}

Ao salvar a configuração, o sistema calcula a razão de contraste da WCAG 2.1 (luminância relativa do sRGB, razão $(L_1 + 0{,}05)/(L_2 + 0{,}05)$) entre cada cor do tema e o branco, conforme a Tabela~\ref{tab:pares-contraste}, e emite um aviso, sem bloquear o salvamento, quando a razão é inferior ao mínimo do par (RN09, \req{RNF05}). A razão é simétrica, e o par é sempre contra o branco: o par ``primária e branco'' cobre a cor primária usada como texto sobre o fundo da página (por exemplo, em links). O texto \emph{sobre} um fundo da cor primária, ao contrário, não é branco fixo: é calculado, como descrito adiante.

\textbf{Por modo (M18; \adr{ADR-025}).} A tabela abaixo descreve o que está integrado (um modo, contra o branco). No alvo, os mesmos três pares, com os mesmos mínimos, são verificados \emph{por modo}: a paleta clara contra \texttt{\#FFFFFF} e a noturna contra o fundo noturno, \texttt{\#0A0A0A} (Tabela~\ref{tab:neutros}). As colunas das cores noturnas são \texttt{darkPrimaryColor}, \texttt{darkSecondaryColor} e \texttt{darkAccentColor}; \texttt{contrastWarnings(registro, modo)} (integrado no M-017; modo claro por omissão) recebe o registro gravado e lê só os três campos do modo, e devolve, em cada aviso, a forma de sempre, \texttt{\{ code, field, ratio, minimum, message \}}: o modo está no próprio \texttt{field} (os campos da paleta noturna começam por \texttt{dark}) e na mensagem, sem uma chave \texttt{mode}, desvio aceito porque o contrato do Core fixa as cinco chaves; \texttt{textColorsOf(paleta, modo)} devolve o texto sobre a primária e sobre a secundária de cada modo (\texttt{readableTextOn} continua função só da cor); o \texttt{PUT} da aparência (do site e do usuário) calcula os dois conjuntos. Os avisos continuam sem bloquear o salvamento (RN09) e, no alvo, também geram a notificação de contraste do dono (RN26), uma por modo, com a chave de deduplicação \texttt{contrast:\{modo\}} na aparência do usuário e \texttt{contrast-site:\{modo\}} na configuração do site (o administrador tem os dois registros, e as chaves não podem colidir; decisão de 05/10/2026), reaberta a cada ocorrência e resolvida quando a paleta passa a atender.

\begin{table}[htb]
\caption{Pares de contraste verificados na configuração}\label{tab:pares-contraste}
\small
\begin{tabularx}{\linewidth}{@{}L{2.8cm}L{3.6cm}L{1.7cm}Y@{}}
\toprule
\textbf{Campo avisado} & \textbf{Par (cor e branco)} & \textbf{Mínimo} & \textbf{Uso que a razão protege}\\
\midrule
\texttt{primaryColor} & primária $\times$ \texttt{\#FFFFFF} & 4,5:1 & Texto na cor primária sobre a página (por exemplo, links); o texto sobre um fundo da cor primária é tratado por \texttt{readableTextOn}, e não por este par\\
\addlinespace
\texttt{secondaryColor} & secundária $\times$ \texttt{\#FFFFFF} & 4,5:1 & Texto branco sobre a barra superior e os cabeçalhos\\
\addlinespace
\texttt{accentColor} & destaque $\times$ \texttt{\#FFFFFF} & 3:1 & Marcadores e realce não textual (item ativo e \en{hover} da navegação): contraste de componentes de interface (WCAG 1.4.11); o foco não usa esta cor\\
\bottomrule
\end{tabularx}
\fonte{Elaborado pelos autores, a partir da resposta ao pedido P-002 (item 5) e da revisão do merge M-005.}
\end{table}

As funções de cálculo pertencem à UI e estão em \arq{src/lib/theme.js} (merge M-005), um módulo puro --- não importa o Next.js nem o Prisma e não usa APIs do navegador ---, que o servidor e o navegador compartilham:

\begin{itemize}
  \item \texttt{contrastRatio(a, b)} devolve a razão da WCAG 2.1 entre duas cores \texttt{\#RRGGBB}, sem arredondar e independentemente da ordem dos argumentos;
  \item \texttt{contrastWarnings(\{ primaryColor, secondaryColor, accentColor \})} devolve a lista de avisos \texttt{\{ code, field, ratio, minimum, message \}}, com o código \texttt{LOW\_CONTRAST}, vazia se nenhuma cor falha; a mensagem vem em português, por exemplo ``A cor primária tem contraste de 3,67:1 com o branco, abaixo do mínimo de 4,5:1 (WCAG 2.1 AA).'';
  \item \texttt{readableTextOn(cor)} devolve \texttt{\#FFFFFF} ou o texto escuro padrão, o de maior razão de contraste com a cor dada (no empate, o branco).
\end{itemize}

A comparação com o mínimo usa a razão \emph{sem arredondar}; o campo \texttt{ratio} e a mensagem trazem a razão \emph{truncada} em duas casas, para que 4,499 não apareça como 4,50 e pareça cumprir o mínimo (por isso a primária padrão aparece como 3,67, e não 3,68). O formulário de tema (M11) usa \texttt{contrastWarnings} na pré-visualização local, e o Route Handler de \texttt{PUT /api/configuration} (Core) a chama, depois da escrita, com as cores já gravadas --- inclusive quando o corpo da requisição não traz cores, pois o \texttt{PUT} é parcial --- e devolve a lista no campo \texttt{warnings} da resposta, ao lado de \texttt{data} (Seção~\ref{sec:arq-api}). O campo está sempre presente na resposta do \texttt{PUT} (lista vazia quando nenhuma cor falha), o aviso não bloqueia o salvamento (o status continua \texttt{200}) e o \texttt{GET} não o devolve. Cores fora do formato \texttt{\#RRGGBB} lançam \texttt{TypeError}: não chegam às funções, pois a validação de formato as recusa antes (RN09).

\textbf{Ordem de execução e resiliência do \texttt{PUT}.} O Route Handler executa, nesta ordem: (1) a verificação de sessão e de origem; (2) a leitura e a validação do corpo (\texttt{400} e \texttt{422}, sem gravar nada); (3) a gravação (\texttt{updateConfiguration}); (4) a revalidação do \en{layout} raiz (RN10); (5) o cálculo dos avisos, sobre o registro devolvido pela gravação; e (6) a resposta \texttt{200}, com \texttt{data} e \texttt{warnings}. O cálculo vem por último de propósito: é um resultado derivado e não pode impedir a gravação nem a revalidação, que são os efeitos que importam. Se ele falhar --- o que só ocorreria com uma cor gravada fora do formato \texttt{\#RRGGBB}, que a validação impede ---, a rota responde \texttt{500} \arq{INTERNAL_ERROR}, pela regra geral da Tabela~\ref{tab:dataerror} (``outro erro''), \emph{depois} de a configuração já estar gravada e o \en{layout} já revalidado. A rota não engole o erro devolvendo \texttt{"warnings": []}, pois isso afirmaria que não há avisos sem tê-los calculado. Portanto, um \texttt{500} do \texttt{PUT} não significa que a gravação falhou: o cliente deve reconsultar a configuração com o \texttt{GET} para conferir o estado, e repetir o \texttt{PUT} com o mesmo corpo produz o mesmo estado.

\textbf{Texto sobre a cor primária.} O branco fixo não atende ao nível AA com a paleta padrão: sobre a primária \texttt{\#3b82f6}, ele dá 3,68:1. Por isso o texto sobre um fundo da cor primária (o botão principal e o item ativo da navegação) usa a variável \arq{--vitrine-on-primary}, exposta como o utilitário \texttt{text-on-primary}, cujo valor é \texttt{readableTextOn(primaryColor)}, isto é, o branco ou o preto (\texttt{\#000000}), o de maior razão de contraste (no empate, o branco); com a paleta padrão, o preto, que dá 5,71:1. O envoltório \texttt{.vitrine-tema} define essa variável por modo (\arq{--vitrine-light-on-primary} e \arq{--vitrine-dark-on-primary}) junto com as cores no seu atributo \texttt{style} (o layout raiz, estático, não define cores), e quem sobrescrever uma primária deve sobrescrever também o texto sobre ela; do contrário, uma primária escura ficaria com texto escuro. Desde o M18 o texto sobre a secundária também é calculado por modo (\texttt{textColorsOf}): com a paleta padrão é branco no modo claro e preto no noturno. A decisão, tomada com o Engenheiro Chefe, consta na \adr{ADR-020}; o preto puro foi adotado no merge M-009.

\textbf{Limites conhecidos.} (i) Com o preto puro como texto escuro (merge M-009), toda cor \texttt{\#RRGGBB} tem um texto, o branco ou o preto, com pelo menos 4,58:1 --- o pior caso, verificado por varredura dos cinzas e de uma grade de 4.096 cores ---; antes, com \texttt{\#171717}, cerca de 6\% das cores ficavam abaixo de 4,5:1. O par da primária, portanto, deixa de proteger o texto sobre ela e protege só o texto na cor primária sobre a página. (ii) A primária padrão \emph{anterior} (até o merge M-012) tinha 3,68:1 também sobre a página, e a de destaque padrão, 2,15:1: com aquela paleta, texto na cor primária não atenderia ao nível AA, e marcadores na cor de destaque ficam abaixo de 3:1; o primeiro \texttt{PUT /api/configuration} com essa paleta devolve, portanto, dois avisos (primária e destaque). A \adr{ADR-031} substitui esses padrões pelos da Tabela~\ref{tab:padroes-cor}, que não geram aviso; a troca ocorreu com o M15 (merge M-012: esquema, \en{seed} e migração). Por isso os links do conteúdo não usam a cor primária: o token \arq{--vitrine-link} tem como padrão a cor do texto (4,5:1 com qualquer paleta) e os links são sempre sublinhados, distinguindo-se sem depender da cor (WCAG 1.4.1); quem quiser links coloridos sobrescreve o token com uma cor já verificada contra o fundo, o que seria um \texttt{[PEDIDO]} ao Core, depois da especificação do modo noturno (\adr{ADR-025}).

\vitrinefigura[tipo=diagrama, arquivo=seq-personalizacao-tema, agente=Arquiteto, largura=0.95,
  origem={doc/diagrams/src/seq-personalizacao-tema.puml},
  descricao={Sequência de personalização: formulário, PUT na API, validação do formato hexadecimal, atualização no banco, revalidação do layout, cálculo dos avisos de contraste sobre as cores gravadas e leitura da configuração pelo RootLayout na próxima visita.}]
  {Sequência de personalização visual}{fig:seq-personalizacao-tema}

\subsection{Leitura da Configuração e o \en{Build}}\label{sec:arq-tema-estatico}

\textbf{Integrado no merge M-021 (\adr{ADR-030}); a aparência do dono é da Parte 3b.} O layout raiz é estático e \emph{não lê o banco}; a configuração do site é lida, em execução e pela mesma guarda, só pelo portal \texttt{/} (layout do grupo de rotas \texttt{(portal)}, que não entra na URL e dá o layout só a \texttt{/}), pelo login e pelo painel (que a injetam no envoltório \texttt{vitrine-tema}, Seção~\ref{sec:arq-tema-padrao}), e a aparência do dono é lida pelo layout de \texttt{/\{handle\}}, que também define o título e a descrição no \texttt{generateMetadata}. Sem dado do banco no layout raiz, as quatro regras abaixo continuam valendo para cada ponto de leitura, e a \emph{página 404 global} sai do \en{build} sempre com os padrões da Tabela~\ref{tab:padroes-cor}, sem depender de revalidação: o que o parágrafo ``Páginas estáticas'' trata como efeito de recriar o contêiner deixa de existir para ela.

A leitura da configuração para o \en{layout} raiz, para os metadados e para o painel está em \arq{src/app/_lib/site.js} (pasta privada do Core, merge M-007) e segue quatro regras (\adr{ADR-021}):

\begin{itemize}
  \item \textbf{Nenhuma consulta no \en{build}.} \texttt{getSiteConfiguration()} devolve \texttt{null}, sem consultar o banco, quando \arq{NEXT_PHASE} vale \texttt{phase-production-build}. É o que sustenta a afirmação de que ``nenhuma consulta ocorre no \en{build}'' (Seção~\ref{sec:arq-implantacao}): o estágio de \en{build} da imagem Docker só define \arq{DATABASE_URL}, de modo que o arquivo, ou até o diretório, pode não existir. Sem a guarda, o \texttt{next build} \emph{falha} ao pré-renderizar \arq{/_not-found} (``Cannot open database because the directory does not exist''; com um banco sem tabelas, ``Invalid \texttt{prisma.configuration.upsert()} invocation''), o que foi medido na revisão do merge M-007;
  \item \textbf{Falha tolerada em execução.} Se a consulta lançar (banco indisponível ou sem as tabelas), o erro é registrado no servidor e valem os padrões do \arq{globals.css}: o site nunca cai por causa da configuração;
  \item \textbf{Uma consulta por requisição.} A função usa o \texttt{cache} do React, de modo que o \en{layout}, o \texttt{generateMetadata} e a página do mesmo pedido (portal, login ou painel) compartilham o resultado (fora do servidor do Next.js, como nos testes, cada chamada consulta);
  \item \textbf{Só cores válidas.} \texttt{buildThemeStyle(registro)}, função pura e sem imports do Next.js no mesmo módulo, monta as variáveis das duas paletas e inclui no \texttt{style} somente os valores que casam com \texttt{\^{}\#[0-9a-fA-F]\{6\}\$}. o texto sobre as cores vem de \texttt{textColorsOf}, por modo, e sem nenhuma cor válida o atributo \texttt{style} é omitido; a UI repete a defesa em \texttt{safeThemeStyle} (\arq{theme.js}), que deixa passar só as dez variáveis conhecidas com valor \texttt{\#RRGGBB}. A verificação é necessária porque o valor lido do banco vai para um atributo HTML e porque \texttt{readableTextOn} lança \texttt{TypeError} com cor inválida. A função fica no Core, e não em \arq{src/lib/theme.js}, por decisão do merge M-007: é pura, e a UI pode movê-la se preferir.
\end{itemize}

\textbf{Páginas estáticas.} O \en{build} não consulta o banco; por isso as páginas pré-renderizadas (o portal \texttt{/} e a página \arq{/_not-found}; as páginas de \texttt{/\{handle\}} nascem sob demanda, sem HTML por \en{handle} no \en{build}) saem dele \emph{sem} o atributo \texttt{style} e com o título padrão (``Vitrine''), e valem os padrões do \arq{globals.css}. O portal passa a refletir a configuração do banco na \emph{primeira revalidação}: a primeira mutação pela API que o revalida (alterar a configuração ou um evento que muda o diretório; Tabela~\ref{tab:revalidacao}) ou o vencimento do prazo; a 404 global fica sempre com os padrões. As rotas dinâmicas (\arq{/login} e \arq{/admin}) já nascem com a configuração. Foi verificado por HTTP, sobre o \en{build} de produção, em duas rodadas independentes (a do Core e a do Arquiteto): antes de qualquer revalidação, \arq{/} e a página 404 respondem \texttt{HIT} com os padrões do CSS; depois de um \texttt{PUT /api/configuration}, a primeira requisição ao portal responde \texttt{MISS} (regeneração) e a seguinte, \texttt{HIT}, ambas com a nova paleta e o novo nome, sem novo \en{build} (o e2e do tema, M-021, repete a medição).

\textbf{Consequência para a implantação (comportamento por projeto).} As páginas regeneradas em execução ficam no sistema de arquivos do contêiner (\arq{.next/server/app/}), e não no volume do banco. Foi medido em contêiner, com a mesma imagem, pelo Core (merge M-008) e, de forma independente, pelo Arquiteto: depois de uma revalidação, um \texttt{docker compose restart} \emph{mantém} as páginas com o tema do banco; já a \emph{recriação} do contêiner --- \texttt{docker compose down} seguido de \texttt{up -d}, uma nova imagem ou \texttt{up -d -{}-force-recreate} --- descarta esse sistema de arquivos, e a página inicial e a 404 voltam a responder \texttt{HIT} com os padrões do CSS e o título ``Vitrine'', enquanto \arq{/login} e \arq{/admin}, dinâmicas, já nascem com a paleta e o nome do banco; a sessão do administrador e a configuração permanecem no volume. Salvar a configuração --- um \texttt{PUT /api/configuration} autenticado, mesmo sem alterar nada, que no painel será o salvamento do formulário de configurações (M11, ainda inexistente) --- revalida o portal: a primeira requisição seguinte responde \texttt{MISS}, já com a paleta e o nome do banco, e as demais, \texttt{HIT}. Isso é \textbf{comportamento por projeto (\en{by design})} da \adr{ADR-021}, e não um defeito: o \en{build} não tem banco e, portanto, não pode embutir o tema nas páginas estáticas. Um passo automático depois do \texttt{up -d} exigiria as credenciais do administrador na implantação ou uma rota de revalidação nova e protegida, que o contrato da API (\adr{ADR-018}) não prevê; por isso o passo é manual e consta do Apêndice~\ref{ap:implantacao} (``Atualização''). Para a 404 e a página inicial do \en{scaffold}, o efeito é só visual. Para as páginas públicas (M07), a regra é \emph{não pré-renderizá-las com dados no \en{build}}: \texttt{generateStaticParams} devolve apenas o que o banco informar (nada, em um \en{build} sem banco) e as páginas nascem sob demanda, na primeira visita, com os dados e o tema do banco em execução; o \texttt{revalidate} por tempo nas páginas públicas (M07) é a alternativa para o que restar estático.

\subsection{Aparência Efetiva: Duas Camadas e Cadeia de Padrões (Épico v1)}\label{sec:arq-tema-efetiva}

\textbf{Situação: integrado (A-025; a camada de dados, A-026, merge M-039; a interface, A-027, merge M-038; a mescla na moldura, no painel e no seletor, A-031, merge M-041; o logotipo, o \en{favicon} e o avatar, A-029, merge M-042; os dois últimos em 09/10/2026; \adr{ADR-038}).} Antes da A-031, o tema de uma página era o da \emph{única} linha de aparência que a rota lia: a do site, no portal, no \texttt{/login} e no painel, e a do dono, no espaço dele (\adr{ADR-030}). Agora a aparência tem duas camadas, e o tema de uma página é o resultado de uma mescla: no espaço, \texttt{getSpace} usa \texttt{getEffectiveAppearance} na geração da página; no painel, \texttt{panelAppearance} (\arq{src/app/admin/\_lib/appearance.js}) devolve o registro do sistema para o administrador e a aparência efetiva para o usuário comum, a cada requisição (o painel é dinâmico).

\textbf{O que vale em cada lugar.}
\begin{itemize}
  \item \emph{Portal} \texttt{/}, \texttt{/login} e \emph{painel do administrador}: só a aparência do \textbf{sistema} (\texttt{Configuration}). Mudar a cor do sistema muda o painel do administrador na hora, e é o que corrige o defeito do item 6 do épico.
  \item \emph{Espaço} de um usuário (inclusive o do administrador) e \emph{painel do usuário comum}: a aparência efetiva, isto é, a \textbf{página do usuário} (\texttt{Appearance}) \emph{sobre} a do sistema.
\end{itemize}

\textbf{A mescla.} Por unidade de herança (Seção~\ref{sec:arq-dados-v1}): cada \emph{paleta de modo} (as três cores juntas) é a do usuário se estiver definida, senão a do sistema (os dois modos são independentes); o \emph{nome}, o \emph{logotipo} e o \en{favicon} seguem a cadeia \textbf{embutido $\rightarrow$ sistema $\rightarrow$ usuário}: vale o do usuário, senão o do sistema, senão o embutido. \texttt{showLogo}, \texttt{showName}, a descrição e o rodapé são só do usuário. O Dados entrega a mescla de dados (\texttt{getEffectiveAppearance}, sem o embutido); o Core aplica o embutido e monta o tema (\texttt{buildThemeStyle}) com a paleta efetiva; a interface recebe tudo por \emph{props}.

\textbf{O embutido.} É o que existe sem configuração alguma, e não ocupa o banco: o logotipo genérico provisório (o componente \texttt{DefaultLogo}, SVG \en{inline} que usa as cores do tema), o \en{favicon} pequeno e neutro (\arq{public/favicon.ico}, 1,1~KiB, a letra ``V'' branca em um quadrado cinza-azulado; o primeiro, de 26~KiB, pesava no LCP, A-024; declarado em \texttt{icons} do layout raiz; a convenção de arquivo de \arq{src/app} o somaria ao favicon enviado em todas as páginas, e o navegador escolheria entre os dois) e o texto \texttt{Vitrine}. Por isso o \en{seed} e os testes não precisam de imagem. Remover um logotipo ou um \en{favicon} enviado volta ao nível de baixo, e o campo nunca fica vazio; a regra da barra superior é só ``ao menos um de \texttt{showLogo} e \texttt{showName} ligado'' (RN31).

\textbf{Contraste.} O aviso e a notificação \texttt{CONTRAST} (RN26) são calculados sobre as paletas que o \emph{dono definiu}; a paleta herdada é responsabilidade de quem a definiu (o administrador, no sistema). Anular a paleta de um modo ``usar a aparência do sistema'' resolve a notificação daquele modo.

\textbf{Estáticas e revalidação.} O layout raiz continua estático e sem dado do banco (\adr{ADR-030}); a mescla é feita na geração do espaço (e a leitura do site segue a guarda do \texttt{getSiteConfiguration}, sem consultar o banco no \en{build}, \adr{ADR-021}). Quando o sistema muda uma unidade que os usuários herdam (a paleta de um modo, o nome, o logotipo ou o \en{favicon}), os espaços que herdam aquela unidade são revalidados por caminhos literais, além do portal, e o par de ícones do botão de modo revalida o espaço de todos (custo medido na Seção~\ref{sec:arq-isr}; o logotipo e o \en{favicon} do sistema revalidam o portal e os espaços que os herdam, e os do usuário, só o espaço dele). Uma mudança na aparência de um usuário revalida só o espaço dele. A barra superior (\texttt{showLogo} e \texttt{showName}), o par de ícones e o título do diretório chegam à moldura do espaço, ao painel, ao portal e ao \texttt{/login} por \emph{props}; \texttt{showLogo} e \texttt{showName} são sempre os do próprio usuário, e o nome do espaço, quando o usuário não definiu um, é o do sistema.

\textbf{O \en{favicon} no \texttt{<head>}.} O \texttt{generateMetadata} do portal, do \texttt{/login}, do painel e do espaço emite \texttt{icons} com o \en{favicon} efetivo (\texttt{/media/<id>}, PNG de 64$\times$64) quando há um enviado; o de um segmento filho substitui o do pai, de modo que o espaço de quem enviou o seu emite \emph{só} o dele, e quem não tem nenhum, nem o sistema, herda o de reserva (\arq{/favicon.ico}). O logotipo chega ao cabeçalho como \texttt{logoUrl} (\texttt{/media/<id>}, ou o embutido).

\textbf{O seletor do administrador.} Em \arq{/admin/aparencia}, o administrador escolhe por \texttt{?alvo=sistema} (o padrão dele; edita \texttt{Configuration} por \texttt{PUT /api/configuration}) ou \texttt{?alvo=pagina} (edita a própria \texttt{Appearance} por \texttt{PUT /api/appearance}); o usuário comum só edita a página. O mesmo formulário compartilhado serve aos dois alvos, e o botão ``restaurar aos padrões'' da página vira ``usar a aparência do sistema''. Depois de salvar, o formulário chama \texttt{router.refresh()}, porque o layout do painel persiste entre navegações do cliente.
